% Source: https://artofproblemsolving.com/community/c3h2411379_geometry_algebra_and_more_problems
% Stable identity: aops-topic-2411379-q11
The vertices of a regular octagon lie on a circle and divide it into $8$ equal arcs of $45^\circ$. Three vertices cut the circle into three arcs containing $a$, $b$ and $c$ of these small arcs, where $a,b,c$ are positive integers with $a+b+c=8$.

The angle of the triangle at a vertex is an inscribed angle standing on the opposite arc, so it equals half of that arc: the angles are $22.5^\circ a$, $22.5^\circ b$ and $22.5^\circ c$. All these triangles have the same circumcircle, so two of them are congruent exactly when they have the same angles, that is, the same multiset $\{a,b,c\}$. (Conversely, any such multiset is realised by vertices $0$, $a$ and $a+b$, counted around the octagon.)

So we count the ways to write $8$ as a sum of three positive integers, ignoring order:
\[8=1+1+6=1+2+5=1+3+4=2+2+4=2+3+3.\]

\begin{asy}
usepackage("amsmath");
size(12cm);
defaultpen(fontsize(10pt));
int[][] P={{1,1,6},{1,2,5},{1,3,4},{2,2,4},{2,3,3}};
for(int t=0;t<5;++t){
  pair c=(2.6*t,0);
  pair V(int k){return c+dir(90-45*k);}
  path oc=V(0); for(int k=1;k<8;++k) oc=oc--V(k); oc=oc--cycle;
  draw(oc,gray(0.6));
  for(int k=0;k<8;++k) dot(V(k),gray(0.4)+linewidth(2));
  int i=0, j=P[t][0], m=P[t][0]+P[t][1];
  fill(V(i)--V(j)--V(m)--cycle,rgb(0.8,0.88,1));
  draw(V(i)--V(j)--V(m)--cycle,blue+linewidth(0.9));
  label("$"+string(P[t][0])+","+string(P[t][1])+","+string(P[t][2])+"$",c+(0,-1.25),S);
}
\end{asy}

There are $\boxed{5}$ non-congruent triangles (shown with their arc counts).
